\specsection{序列空间上的测度}

经典概率论发展最充分的分支之一是极限定理的分支（大数定律、趋于 Gauss--Laplace 律，\ldots)；这涉及对由牵涉极大总体的现象所显示的统计规律性概念的深化。这些问题的正确数学表述需要在序列空间上引入测度；这些空间构成有限维空间最明显的推广，是 Fréchet、Lévy、Lusin\ldots{} 在 1920 年前后所从事的“广义分析”研究的首选对象。新概率论方法的创立者 Khintchine 与 Kolmogoroff 都是 Lusin 的学生，而 Lévy 很快就转向了概率问题，这也并非偶然：这些构成了新方法的试金石。

序列空间上的测度第一次隐含地出现于 E. Borel 1909 年讨论可数概率的著作 (IV) 中。Borel 的一个高度独创的想法，是把他刚得到的概率结果应用于证明几乎每个介于 0 与 1 之间的实数的十进制展开所具有的性质。这一应用基于下列基本注记：用给定基 \(q\)（\(q \geqslant 2\)）下展开式中的数字（或“数码”）的序列来定义每个介于 0 与 1 之间的实数；如果依次随机地抽取一个数 \(x\) 的各个数字，以相互独立、相同的概率 \(1 / q\)（对 \(0,1,\ldots ,q - 1\)），则 \(x\) 落在 {[}0,1{[} 的一个区间中的概率等于该区间的长度。

1923 年，Steinhaus (V) 严格地建立了这些结果，并描述了 Borel 所考虑的无穷随机抽取序列的精确数学模型：为简单起见取 q=2，并用 I 表示两个元素的集合 \(\{0,1\}\)；给 I 装备测度 \(\mu\)，它由 \(\mu(0)=\mu(1)=\frac{1}{2}\) 定义；乘积空间 \(I^{N}\) 的元素是由等于 0 或 1 的数组成的序列 \(\varepsilon=(\varepsilon(n))_{n\in\mathbf{N}}\)，而映射 \(\varphi:\varepsilon\mapsto\sum_{n\geqslant0}\varepsilon(n)\cdot2^{-n-1}\) 除一个可数集外是 \(I^{N}\) 到区间 {[}0,1{]} 上的一个双射；此外，\(\varphi^{-1}\) 把 {[}0,1{]} 上的 Lebesgue 测度变换为 \(I^{N}\) 上的测度 P，即各因子上测度 \(\mu\) 的乘积。实际上，Steinhaus 当时并没有乘积测度的构造可用；他利用拟双射 \(\varphi\) 的存在性，构造 \(I^{N}\) 上的测度 P（从 {[}0,1{]} 上的 Lebesgue 测度出发），然后给出 P 的一个公理化刻画。这样得到的同构使人们能够把概率的语言翻译为测度的语言，并应用关于 Lebesgue 积分的已知定理。

在同一著作中，Steinhaus 考虑了随机级数 \(\sum_{n\geqslant 0}\sigma_n\cdot a_n\)，其中符号 \(\sigma_{n} = \pm 1\) 相互独立且以相同的概率 \(\frac{1}{2}\) 随机选取；在 1928 年至 1935 年之间，他研究了许多其他的随机级数。另一方面，Paley、Wiener 与 Zygmund 考虑了形如 \(\sum_{n = -\infty}^{\infty}a_n\exp (2\pi i(nt + \Phi_n))\) 的随机 Fourier 级数(1)；其中“振幅”\(a_{n}\) 是固定的，而“相位”\(\Phi_{n}\) 是在 {[}0,1{]} 上均匀分布的独立随机变量。虽然从一个问题到另一个问题解析上的困难相差悬殊，但用测度论语言所作的翻译在所有情形中都是相同的，并且是 Borel 与 Steinhaus 所处理情形的一个推广；问题在于构造 \(\mathbf{R}^{\mathbf{N}}\) 上的一个测度，它是一个测度族的乘积，族中的测度都等同于同一个质量为 1 的正测度 \(\mu\)（\(\mathbf{R}\) 上的）；例如，上述随机 Fourier 级数对应于 \(\mu\) 是 {[}0,1{]} 上的 Lebesgue 测度的情形。

为构造这样的乘积测度，可以用两种方法。第一种是直接方法，由 Daniell (VI, a)) 于 1918 年首次精确阐述；它在 1934 年被 Jessen (VII) 重新发现，后者对 \(\mu\) 是 {[}0, 1{]} 上 Lebesgue 测度的情形作了详细研究。第二种方法是寻求类似于 Steinhaus 的技巧，以化归到 {[}0, 1{]} 上的 Lebesgue 测度；在还没有一般测度论的完整阐述可供使用时，这种做法有方便的优点，因为它允许使用 Lebesgue 的定理而无需新的证明。\(^{(2)}\)

